% ============================================================
% Chapter 1: Foundations of Information Security and the CIA Triad
% Language: Greek — edit freely; regenerate from src/content if preferred.
% ============================================================
\chapter{Θεμελιώδεις Έννοιες της Ασφάλειας Πληροφοριών και η Τριάδα CIA}\label{ch:1}
\chaptermeta{Ορισμοί · Η τριάδα CIA και οι επεκτάσεις της · Αρχές σχεδίασης · Ιστορική εξέλιξη · Μοντέλα ασφάλειας · Κριτήρια αξιολόγησης}{Επίπεδο: Εισαγωγικό επίπεδο · Έμφαση στη θεωρία}{Χρόνος μελέτης: 6–8 ώρες}

Κάθε επόμενο κεφάλαιο του βιβλίου —είτε αφορά λειτουργικά συστήματα, κρυπτογραφία, διαχείριση ταυτότητας είτε απόκριση σε περιστατικά— στηρίζεται σε ένα μικρό σύνολο κοινών εννοιών. Το παρόν κεφάλαιο εισάγει αυτό το λεξιλόγιο. Εξηγεί τι εννοούμε όταν μιλάμε για \emph{ασφάλεια}, ποιες ιδιότητες της πληροφορίας επιδιώκουμε να προστατεύσουμε και ποιες αρχές σχεδίασης εφαρμόζουν οι έμπειροι μηχανικοί πριν επιλέξουν οποιοδήποτε συγκεκριμένο προϊόν ή τεχνολογία.

Η παρουσίαση κινείται από τους ορισμούς στις αρχές και, στη συνέχεια, στα τυπικά μοντέλα. Αρχικά διακρίνουμε την κυβερνοασφάλεια από τις ευρύτερες έννοιες της διασφάλισης πληροφοριών και της κυβερνοανθεκτικότητας. Ακολούθως εξετάζουμε την τριάδα CIA και τις ιδιότητες που τη συμπληρώνουν, τις κλασικές αρχές σχεδίασης των Saltzer και Schroeder, καθώς και τα ιστορικά γεγονότα που διαμόρφωσαν τη σύγχρονη αμυντική σκέψη. Το κεφάλαιο ολοκληρώνεται με τα τυπικά μοντέλα (Bell–LaPadula, Biba, Clark–Wilson) και τα σχήματα αξιολόγησης που επιτρέπουν να διατυπώνονται οι ισχυρισμοί ασφάλειας με ακρίβεια και να επαληθεύονται ανεξάρτητα.

\begin{outcomesbox}{Μαθησιακά αποτελέσματα}
Μετά τη μελέτη του κεφαλαίου θα πρέπει να μπορείτε:
\begin{itemize}
  \item Να ορίζετε την κυβερνοασφάλεια, τη διασφάλιση πληροφοριών και την κυβερνοανθεκτικότητα και να εξηγείτε σε τι διαφέρει το εύρος καθεμιάς.
  \item Να αναλύετε μια απαίτηση ασφάλειας σε εμπιστευτικότητα, ακεραιότητα και διαθεσιμότητα και να αναγνωρίζετε πότε απαιτούνται πρόσθετες ιδιότητες (αυθεντικότητα, λογοδοσία, μη αποποίηση, κατοχή, χρησιμότητα).
  \item Να εφαρμόζετε τις αρχές του ελάχιστου προνομίου, της άμυνας σε βάθος, των ασφαλών προεπιλογών και της ασφάλειας εκ σχεδιασμού κατά την αξιολόγηση μιας αρχιτεκτονικής.
  \item Να διακρίνετε την απειλή, την ευπάθεια και τον κίνδυνο και να περιγράφετε πώς τα μέτρα ελέγχου μειώνουν την πιθανότητα ή τις συνέπειες.
  \item Να εξηγείτε τον σκοπό των τυπικών μοντέλων ασφάλειας και των ανεξάρτητων σχημάτων αξιολόγησης, όπως τα Common Criteria.
\end{itemize}
\end{outcomesbox}

\section{Ορισμοί: κυβερνοασφάλεια, διασφάλιση πληροφοριών και κυβερνοανθεκτικότητα}\label{sec:1.1}
Η \textbf{κυβερνοασφάλεια} είναι ο επιστημονικός και επαγγελματικός κλάδος που προστατεύει δικτυωμένα συστήματα, δεδομένα και υπηρεσίες από μη εξουσιοδοτημένη πρόσβαση, τροποποίηση, διακοπή και καταστροφή, διασφαλίζοντας ταυτόχρονα τη νόμιμη λειτουργία τους ακόμη και υπό εχθρικές συνθήκες. Η φράση \emph{υπό εχθρικές συνθήκες} έχει ιδιαίτερη σημασία: σε αντίθεση με τη μηχανική αξιοπιστίας, η οποία αντιμετωπίζει τυχαία σφάλματα, η ασφάλεια προϋποθέτει έναν ευφυή αντίπαλο που αναζητά συστηματικά το πιο αδύναμο σημείο.

Η \textbf{διασφάλιση πληροφοριών} (Information Assurance, IA) έχει ευρύτερο πεδίο. Συνδυάζει τεχνικά, διαδικαστικά και διοικητικά μέτρα, ώστε να διατηρείται \emph{τεκμηριωμένη εμπιστοσύνη} στην πληροφορία σε όλο τον κύκλο ζωής της —από τη δημιουργία και την αποθήκευση έως τη μετάδοση, τη χρήση και την τελική καταστροφή της. Επομένως, η διασφάλιση πληροφοριών δεν αφορά μόνο την αντίσταση σε επιθέσεις, αλλά και τη γνησιότητα, την προέλευση και τη διαρκή χρησιμότητα της πληροφορίας.

Η \textbf{κυβερνοανθεκτικότητα} διευρύνει και τις δύο έννοιες. Περιγράφει την ικανότητα ενός οργανισμού να προβλέπει, να αντέχει, να ανακάμπτει και να προσαρμόζεται σε δυσμενή κυβερνογεγονότα, συμπεριλαμβανομένων εκείνων που επιτυγχάνουν παρά τα προληπτικά μέτρα. Με απλά λόγια, η κυβερνοασφάλεια ρωτά \emph{πώς αποτρέπουμε και εντοπίζουμε μια παραβίαση;}, η διασφάλιση πληροφοριών \emph{μπορούμε να εμπιστευτούμε την πληροφορία μας;} και η ανθεκτικότητα \emph{μπορούμε να συνεχίσουμε να λειτουργούμε και πόσο γρήγορα θα ανακάμψουμε;}

\begin{table}[htbp]
  \centering\small
  \caption{Τρεις συγγενείς αλλά διακριτές έννοιες}
  \begin{tabularx}{\linewidth}{>{\raggedright\arraybackslash}X >{\raggedright\arraybackslash}X >{\raggedright\arraybackslash}X >{\raggedright\arraybackslash}X}
    \toprule
    \textbf{Έννοια} & \textbf{Βασικό ερώτημα} & \textbf{Τυπικά μέτρα} & \textbf{Συνέπεια αν παραμεληθεί} \\
    \midrule
    Κυβερνοασφάλεια & Μπορεί ένας αντίπαλος να παραβιάσει ή να διαταράξει το σύστημα; & Τείχη προστασίας, IDS/IPS, ενημερώσεις, MFA, κρυπτογράφηση & Εισβολή, διαρροή δεδομένων, διακοπή υπηρεσίας \\
    Διασφάλιση πληροφοριών & Μπορούμε να εμπιστευτούμε την πληροφορία; & Κατακερματισμός, ψηφιακές υπογραφές, προέλευση, διαβάθμιση & Αθόρυβη αλλοίωση, αποποίηση ευθύνης, παραπληροφόρηση \\
    Κυβερνοανθεκτικότητα & Μπορούμε να λειτουργούμε παρά την αστοχία; & Αντίγραφα ασφαλείας, εφεδρεία, σχέδια απόκρισης, BCP/DR & Παρατεταμένη διακοπή, μη ανακτήσιμη απώλεια \\
    \bottomrule
  \end{tabularx}
\end{table}
\section{Η τριάδα CIA και οι επεκτάσεις της}\label{sec:1.2}
Η \textbf{τριάδα CIA} —Εμπιστευτικότητα (Confidentiality), Ακεραιότητα (Integrity) και Διαθεσιμότητα (Availability)— αποτελεί τον καθιερωμένο τρόπο ανάλυσης μιας απαίτησης ασφάλειας στους βασικούς της στόχους. Η \textbf{εμπιστευτικότητα} περιορίζει την αποκάλυψη της πληροφορίας μόνο σε εξουσιοδοτημένα μέρη· επιβάλλεται συνήθως μέσω ελέγχου πρόσβασης και κρυπτογράφησης. Η \textbf{ακεραιότητα} προστατεύει την πληροφορία από μη εξουσιοδοτημένη ή μη ανιχνεύσιμη τροποποίηση, ώστε τα δεδομένα να παραμένουν ακριβή και πλήρη· σε αυτόν τον στόχο συμβάλλουν οι συναρτήσεις κατακερματισμού, οι ψηφιακές υπογραφές, ο έλεγχος εκδόσεων και οι επικυρωμένες συναλλαγές. Η \textbf{διαθεσιμότητα} εξασφαλίζει έγκαιρη και αξιόπιστη πρόσβαση στους εξουσιοδοτημένους χρήστες και υποστηρίζεται από πλεονασμό, σχεδιασμό χωρητικότητας, περιορισμό ρυθμού αιτημάτων και διαδικασίες ανάκαμψης.

Ορισμένες πρόσθετες ιδιότητες εξειδικεύουν την τριάδα χωρίς να την αντικαθιστούν. Η \textbf{αυθεντικότητα} βεβαιώνει ότι ένα δεδομένο ή ένα μέρος είναι πράγματι αυτό που ισχυρίζεται ότι είναι. Η \textbf{λογοδοσία} συνδέει κάθε ενέργεια με ένα αναγνωρίσιμο υποκείμενο μέσω αξιόπιστων αρχείων καταγραφής. Η \textbf{μη αποποίηση} παρέχει αποδεικτικά στοιχεία που καθιστούν δύσκολη τη μεταγενέστερη άρνηση μιας ενέργειας· η ισχύς της, ωστόσο, εξαρτάται τελικά από τη φύλαξη των κλειδιών, την ταυτοποίηση των προσώπων και το νομικό πλαίσιο.

Οι τρεις στόχοι αλληλεπιδρούν και συχνά ανταγωνίζονται μεταξύ τους. Η κρυπτογράφηση ενισχύει την εμπιστευτικότητα, όμως αν χαθούν τα κλειδιά τα δεδομένα καθίστανται μη διαθέσιμα. Ο πλεονασμός βελτιώνει τη διαθεσιμότητα, αλλά κάθε επιπλέον αντίγραφο είναι ένα ακόμη σύστημα που πρέπει να προστατεύεται και να διατηρείται συνεπές. Μια τεκμηριωμένη αρχιτεκτονική δηλώνει, επομένως, ποιους στόχους υπηρετεί κάθε μέτρο, εντοπίζει τυχόν νέα έκθεση που αυτό εισάγει και εξηγεί πώς θα αντιμετωπιστεί ο κίνδυνος που απομένει (υπολειπόμενος κίνδυνος).

\begin{examplebox}{Παράδειγμα εφαρμογής}
Ένα νοσοκομείο κρυπτογραφεί τη βάση δεδομένων ασθενών (εμπιστευτικότητα) και τη συγχρονίζει με δεύτερη εγκατάσταση (διαθεσιμότητα). Αν το μοναδικό αντίγραφο του κλειδιού αποκρυπτογράφησης βρίσκεται στον κύριο διακομιστή και αυτός καταστραφεί από λυτρισμικό, το αντίγραφο της βάσης είναι άχρηστο: η εμπιστευτικότητα διατηρήθηκε εις βάρος της διαθεσιμότητας. Το μέτρο που συμφιλιώνει τους δύο στόχους είναι η ορθή διαχείριση κλειδιών.
\end{examplebox}

\section{Αρχές σχεδίασης που διέπουν όλα τα επόμενα κεφάλαια}\label{sec:1.3}
Τέσσερις αρχές επανέρχονται σε όλο το βιβλίο ως αρχιτεκτονικοί περιορισμοί. Η \textbf{ασφάλεια εκ σχεδιασμού} (security by design) απαιτεί οι ιδιότητες ασφάλειας να προσδιορίζονται, να μοντελοποιούνται, να υλοποιούνται και να επαληθεύονται από την αρχή του κύκλου ζωής και όχι να προστίθενται μετά την παραγωγική λειτουργία, όταν οι αλλαγές είναι δαπανηρές και αποσπασματικές. Η \textbf{άμυνα σε βάθος} συνδυάζει ανεξάρτητα και ποικίλα προληπτικά, ανιχνευτικά και διορθωτικά μέτρα, ώστε η αστοχία ενός μεμονωμένου επιπέδου να μην κρίνει την έκβαση. Η \textbf{αρχή του ελάχιστου προνομίου} παραχωρεί σε κάθε χρήστη, διεργασία ή υπηρεσία μόνο την εξουσία που απαιτείται για συγκεκριμένη εργασία και για συγκεκριμένο χρονικό διάστημα, και την ανακαλεί όταν δεν είναι πλέον αναγκαία. Οι \textbf{ασφαλείς προεπιλογές} (fail-safe defaults) απορρίπτουν την πρόσβαση εκτός αν έχει ρητά επιτραπεί και απαιτούν από τα συστατικά ενός συστήματος να καταλήγουν, σε περίπτωση αστοχίας, σε γνωστή ασφαλή κατάσταση.

Οι αρχές αυτές προέρχονται από το θεμελιώδες άρθρο των Saltzer και Schroeder (1975), το οποίο εισήγαγε επίσης την \emph{οικονομία μηχανισμού} (απλοί και μικροί σχεδιασμοί, ώστε να μπορούν να ελεγχθούν), την \emph{πλήρη διαμεσολάβηση} (έλεγχος κάθε πρόσβασης και όχι μόνο της πρώτης), τον \emph{ανοιχτό σχεδιασμό} (η ασφάλεια δεν πρέπει να εξαρτάται από τη μυστικότητα του σχεδιασμού), τον \emph{διαχωρισμό προνομίων} (περισσότερες από μία ανεξάρτητες προϋποθέσεις για ευαίσθητες ενέργειες), τον \emph{ελάχιστο κοινό μηχανισμό} (περιορισμός των κοινόχρηστων στοιχείων μεταξύ χρηστών) και την \emph{ψυχολογική αποδοχή} (η ασφάλεια πρέπει να είναι εύχρηστη, διαφορετικά οι χρήστες θα την παρακάμψουν).

Οι αρχές είναι συμπληρωματικές και πρέπει να εφαρμόζονται από κοινού. Η πολυεπίπεδη άμυνα χωρίς ελάχιστο προνόμιο απλώς πολλαπλασιάζει τους ισχυρούς λογαριασμούς που μπορεί να υποκλέψει ένας επιτιθέμενος· μια αυστηρή πολιτική άρνησης εξ ορισμού χωρίς λειτουργικές διαδικασίες ωθεί τους χρήστες σε μη ασφαλείς παρακάμψεις. Μια χρήσιμη επαγγελματική πρακτική είναι να κατονομάζεται η αρχή που παραβιάστηκε σε κάθε ανασκόπηση μετά από περιστατικό· για παράδειγμα, μια τραπεζική μεταφορά που εγκρίνεται από ένα μόνο πρόσωπο παραβιάζει τον διαχωρισμό προνομίων.

\section{Ιστορική εξέλιξη: από το Creeper στον υπερσυνδεδεμένο οργανισμό}\label{sec:1.4}
Η σύγχρονη αρχιτεκτονική ασφάλειας γίνεται καλύτερα κατανοητή ως συσσώρευση διδαγμάτων και όχι ως αλληλουχία απειλών που η μία αντικατέστησε την άλλη. Τα πρώιμα πειράματα στο ARPANET, όπως το \emph{Creeper} και το αντίστοιχο \emph{Reaper} (αρχές της δεκαετίας του 1970), έδειξαν τόσο τη δυνατότητα αυτοδιαδιδόμενου κώδικα όσο και την αυτοματοποιημένη αφαίρεσή του. Το \emph{Morris Worm} του 1988 ανέδειξε πώς ένα μεμονωμένο σφάλμα λογισμικού μπορεί να εξαπλωθεί αλυσιδωτά σε διασυνδεδεμένους υπολογιστές και οδήγησε άμεσα στη δημιουργία της πρώτης Ομάδας Απόκρισης σε Υπολογιστικά Περιστατικά (CERT).

Στη συνέχεια, το εμπορικό Διαδίκτυο και ο Παγκόσμιος Ιστός διεύρυναν δραματικά την επιφάνεια επίθεσης: οι απομακρυσμένα προσβάσιμες υπηρεσίες, η έγχυση SQL, το cross-site scripting, τα botnet και το οργανωμένο κυβερνοέγκλημα έγιναν καθημερινά ζητήματα. Τις επόμενες δεκαετίες, το υπολογιστικό νέφος, οι κινητές πλατφόρμες, το λυτρισμικό ως υπηρεσία, οι επιθέσεις με επίκεντρο την ταυτότητα, οι παραβιάσεις της εφοδιαστικής αλυσίδας λογισμικού και η σύγκλιση της πληροφορικής με την επιχειρησιακή τεχνολογία (OT) εισήγαγαν νέα διοικητικά πεδία και νέες εξαρτήσεις.

Κάθε εποχή στηρίχθηκε σε μια παραδοχή εμπιστοσύνης που οι επιτιθέμενοι τελικά διέψευσαν. Ο σημερινός οργανισμός είναι υβριδικός, έντονα διασυνδεδεμένος και εξαρτημένος από εξωτερικούς παρόχους ταυτότητας, λογισμικού και υποδομών. Το πρακτικό συμπέρασμα είναι ότι η θέση μέσα στο δίκτυο δεν αρκεί πλέον για να θεμελιώσει εμπιστοσύνη· η ταυτότητα, η κατάσταση της συσκευής και η εξουσιοδότηση πρέπει να επαληθεύονται συνεχώς. Αυτή είναι η βασική ιδέα της αρχιτεκτονικής μηδενικής εμπιστοσύνης (Zero Trust), την οποία εξετάζει αναλυτικά το Κεφάλαιο 13.

\begin{table}[htbp]
  \centering\small
  \caption{Παραδοχές εμπιστοσύνης και οι αστοχίες που τις διέψευσαν}
  \begin{tabularx}{\linewidth}{>{\raggedright\arraybackslash}X >{\raggedright\arraybackslash}X >{\raggedright\arraybackslash}X >{\raggedright\arraybackslash}X}
    \toprule
    \textbf{Περίοδος} & \textbf{Παραδοχή εμπιστοσύνης} & \textbf{Χαρακτηριστική αστοχία} & \textbf{Καινοτομία ελέγχου που προέκυψε} \\
    \midrule
    1970–1980 & Αξιόπιστοι υπολογιστές και χρήστες & Creeper, Morris Worm & Έλεγχος πρόσβασης, CERT \\
    1990–2000 & Ασφαλές εσωτερικό, εχθρικό εξωτερικό & Code Red, SQL Slammer & Τείχη προστασίας, διαχείριση ενημερώσεων, antivirus \\
    2010–2020 & Αξιόπιστοι προμηθευτές, περίμετρος VPN & Παραβίαση Target, WannaCry & Κατάτμηση δικτύου, EDR, MFA \\
    2020 κ.ε. & Ποτέ μην εμπιστεύεσαι, πάντα επαλήθευε & SolarWinds, λυτρισμικό σε OT & Zero Trust, SBOM, XDR \\
    \bottomrule
  \end{tabularx}
\end{table}
\section{Απειλή, ευπάθεια και κίνδυνος: το βασικό λεξιλόγιο}\label{sec:1.5}
Τρεις όροι χρησιμοποιούνται συχνά εναλλακτικά στην καθημερινή γλώσσα, έχουν όμως διακριτό τεχνικό περιεχόμενο. \textbf{Απειλή} είναι κάθε περίσταση, γεγονός ή δράστης που μπορεί να προκαλέσει βλάβη —μια εγκληματική ομάδα, ένας δυσαρεστημένος υπάλληλος, μια πλημμύρα. \textbf{Ευπάθεια} είναι μια αδυναμία που μπορεί να εκμεταλλευτεί ή να ενεργοποιήσει μια απειλή —μια υπηρεσία χωρίς ενημερώσεις, ένας αδύναμος κωδικός πρόσβασης, η απουσία αντιγράφων ασφαλείας. \textbf{Κίνδυνος} είναι η δυνατότητα δυσμενούς επίπτωσης υπό συνθήκες αβεβαιότητας· υπάρχει μόνο όταν μια σχετική απειλή συναντά μια εκμεταλλεύσιμη ευπάθεια σε ένα περιουσιακό στοιχείο που έχει αξία.

Η γνωστή έκφραση \emph{κίνδυνος ≈ πιθανότητα × επίπτωση} είναι χρήσιμο μοντέλο ιεράρχησης και όχι φυσικός νόμος. Κάθε οργανισμός οφείλει να ορίζει τις δικές του κλίμακες, να δηλώνει τις παραδοχές του και να αναγνωρίζει την αβεβαιότητα. \textbf{Μέτρο ελέγχου} είναι κάθε ενέργεια που μεταβάλλει την πιθανότητα, την επίπτωση ή και τα δύο: ο έλεγχος ταυτότητας πολλαπλών παραγόντων μειώνει την πιθανότητα κατάληψης λογαριασμού, ενώ τα αμετάβλητα αντίγραφα ασφαλείας μειώνουν την επίπτωση ενός λυτρισμικού.

Μετά την εφαρμογή των μέτρων απομένει ο \textbf{υπολειπόμενος κίνδυνος}, ο οποίος πρέπει να αντιμετωπίζεται ρητά. Υπάρχουν τέσσερις επιλογές: η περαιτέρω \emph{μείωση}, η \emph{αποδοχή} (από στέλεχος που έχει τη σχετική αρμοδιότητα), η \emph{μεταφορά} όπου αυτό έχει νόημα (για παράδειγμα, μέσω ασφάλισης ή συμβάσεων) και η \emph{αποφυγή}, δηλαδή η εγκατάλειψη της επικίνδυνης δραστηριότητας. Τα ώριμα προγράμματα ασφάλειας αξιολογούν το σύνολο των μέτρων τους με βάση την κάλυψη, την ανεξαρτησία και την πραγματική λειτουργική αποτελεσματικότητα —όχι με βάση τον αριθμό των προϊόντων που έχουν εγκαταστήσει.

\section{Η εξάδα του Parker και ο κύβος του McCumber}\label{sec:1.6}
Η τριάδα CIA είναι αναγκαία, αλλά δεν είναι πάντοτε επαρκής. Ο Donn Parker πρότεινε την \textbf{εξάδα του Parker} (Parkerian Hexad), η οποία προσθέτει στις τρεις κλασικές ιδιότητες την \emph{αυθεντικότητα}, την \emph{κατοχή (έλεγχο)} και τη \emph{χρησιμότητα}. Οι προσθήκες αυτές αποτυπώνουν καταστάσεις που η τριάδα περιγράφει ελλιπώς. Όταν κλαπεί ένας κρυπτογραφημένος φορητός υπολογιστής, η εμπιστευτικότητα μπορεί να παραμένει άθικτη, όμως η κατοχή έχει σαφώς απολεσθεί. Όταν τα κρυπτογραφημένα δεδομένα διασώζονται αλλά τα κλειδιά έχουν χαθεί οριστικά, τα δεδομένα παραμένουν εμπιστευτικά, έχουν όμως χάσει κάθε χρησιμότητα.

Ο \textbf{κύβος του McCumber} προσφέρει μια συμπληρωματική, τρισδιάστατη οπτική. Εξετάζει την ασφάλεια στις τρεις \emph{καταστάσεις} της πληροφορίας (αποθήκευση, επεξεργασία, μετάδοση), στις τρεις \emph{κατηγορίες μέτρων προστασίας} (τεχνολογία, πολιτικές και πρακτικές, άνθρωποι) και στους ίδιους τους στόχους ασφάλειας. Η συνδυαστική χρήση των δύο μοντέλων αποκαλύπτει ελλιπείς ισχυρισμούς όπως «τα δεδομένα είναι κρυπτογραφημένα, άρα είναι ασφαλή». Ο αξιολογητής καλείται να ρωτήσει: ποιες ιδιότητες προστατεύονται, σε ποια κατάσταση της πληροφορίας και με ποιον συνδυασμό τεχνολογίας, διαδικασιών και ανθρώπινης συμπεριφοράς;

\section{Τυπικά μοντέλα ασφάλειας: Bell–LaPadula, Biba και Clark–Wilson}\label{sec:1.7}
Οι αρχές μάς λένε \emph{τι} επιδιώκουμε· τα τυπικά μοντέλα ορίζουν \emph{με ακρίβεια} ποιες συμπεριφορές ενός συστήματος επιτρέπονται, ώστε ένας σχεδιασμός να μπορεί να αναλυθεί και, ιδανικά, να αποδειχθεί ορθός. Το μοντέλο \textbf{Bell–LaPadula (BLP)}, που αναπτύχθηκε για στρατιωτικά συστήματα τη δεκαετία του 1970, προστατεύει την εμπιστευτικότητα με δύο κανόνες: \emph{όχι ανάγνωση προς τα πάνω} (ένα υποκείμενο δεν μπορεί να διαβάσει αντικείμενα υψηλότερης διαβάθμισης) και \emph{όχι εγγραφή προς τα κάτω} (ένα υποκείμενο δεν μπορεί να γράψει πληροφορία σε χαμηλότερη διαβάθμιση, κάτι που θα διέρρεε απόρρητα).

Το μοντέλο \textbf{Biba} αποτελεί το αντίστοιχο του BLP για την ακεραιότητα και αντιστρέφει τους κανόνες του: \emph{όχι ανάγνωση προς τα κάτω} (μια αξιόπιστη διεργασία δεν πρέπει να καταναλώνει λιγότερο αξιόπιστα δεδομένα) και \emph{όχι εγγραφή προς τα πάνω} (ένα μη αξιόπιστο υποκείμενο δεν μπορεί να τροποποιεί πιο αξιόπιστα αντικείμενα). Στα εμπορικά περιβάλλοντα, ωστόσο, ενδιαφέρουν λιγότερο τα επίπεδα διαβάθμισης και περισσότερο η ορθότητα των συναλλαγών. Γι' αυτό το μοντέλο \textbf{Clark–Wilson} επιβάλλει την ακεραιότητα μέσω \emph{καλά σχηματισμένων συναλλαγών}, οι οποίες εκτελούνται μόνο από εξουσιοδοτημένους χρήστες μέσω πιστοποιημένων προγραμμάτων, σε συνδυασμό με τον \emph{διαχωρισμό καθηκόντων} και την πλήρη καταγραφή ελέγχου.

\begin{notebox}{Γιατί έχει πρακτική σημασία}
Τα υποχρεωτικά επίπεδα ακεραιότητας των Windows (Κεφάλαιο 2) προέρχονται άμεσα από το μοντέλο Biba, ενώ η έγκριση τραπεζικών μεταφορών από δύο πρόσωπα —δημιουργό και ελεγκτή (Κεφάλαιο 5)— αποτελεί κλασικό μέτρο τύπου Clark–Wilson.
\end{notebox}

\section{Κριτήρια αξιολόγησης: από το TCSEC στα Common Criteria}\label{sec:1.8}
Πώς μπορεί ένας πελάτης να γνωρίζει αν ένα προϊόν παρέχει πράγματι την ασφάλεια που ισχυρίζεται ο κατασκευαστής του; Τα σχήματα αξιολόγησης απαντούν σε αυτό το ερώτημα μέσω ανεξάρτητης και επαναλήψιμης αξιολόγησης. Τα αμερικανικά \textbf{Trusted Computer System Evaluation Criteria (TCSEC)}, γνωστά ως \emph{Πορτοκαλί Βιβλίο} (Orange Book, 1983), κατέτασσαν τα λειτουργικά συστήματα σε κατηγορίες από D (ελάχιστη προστασία) έως A1 (τυπικά επαληθευμένος σχεδιασμός). Η Ευρώπη ακολούθησε με το \textbf{ITSEC}, το οποίο διαχώρισε τη λειτουργικότητα από τη διασφάλιση.

Και τα δύο αντικαταστάθηκαν από τα διεθνή \textbf{Common Criteria (ISO/IEC 15408)}. Στο πλαίσιο αυτό, ένα \emph{Προφίλ Προστασίας} (Protection Profile) περιγράφει τις απαιτήσεις ασφάλειας για μια κατηγορία προϊόντων (για παράδειγμα, τείχη προστασίας), ενώ ο \emph{Στόχος Ασφάλειας} (Security Target) του κατασκευαστή δηλώνει πώς ένα συγκεκριμένο προϊόν τις ικανοποιεί. Διαπιστευμένα εργαστήρια αξιολογούν στη συνέχεια το προϊόν σε ένα \textbf{Επίπεδο Διασφάλισης Αξιολόγησης (EAL1–EAL7)}. Είναι κρίσιμο να κατανοηθεί τι σημαίνει το EAL: μετρά την \emph{αυστηρότητα} της αξιολόγησης και όχι την \emph{ισχύ} του προϊόντος. Ένα προϊόν EAL4 έχει εξεταστεί διεξοδικότερα από ένα προϊόν EAL2, αλλά μόνο ως προς τις απαιτήσεις που δήλωσε ο δικός του Στόχος Ασφάλειας.

\section*{Βασικοί όροι}
\begin{description}[style=nextline,leftmargin=1.2em]
  \item[Τριάδα CIA] Εμπιστευτικότητα, ακεραιότητα και διαθεσιμότητα: οι τρεις βασικοί στόχοι της ασφάλειας πληροφοριών.
  \item[Μη αποποίηση] Αποδεικτικά στοιχεία που εμποδίζουν ένα μέρος να αρνηθεί αξιόπιστα μια ενέργεια που εκτέλεσε.
  \item[Ελάχιστο προνόμιο] Παραχώρηση μόνο της ελάχιστης αναγκαίας εξουσίας για μια εργασία και για τον ελάχιστο αναγκαίο χρόνο.
  \item[Άμυνα σε βάθος] Διαστρωμάτωση ανεξάρτητων και ποικίλων μέτρων, ώστε καμία μεμονωμένη αστοχία να μην είναι καθοριστική.
  \item[Υπολειπόμενος κίνδυνος] Ο κίνδυνος που απομένει μετά την εφαρμογή των μέτρων και πρέπει να αντιμετωπιστεί ρητά.
  \item[EAL] Επίπεδο Διασφάλισης Αξιολόγησης (Common Criteria): η αυστηρότητα μιας ανεξάρτητης αξιολόγησης προϊόντος.
\end{description}

\section*{Σύνοψη κεφαλαίου}
\begin{itemize}
  \item Η κυβερνοασφάλεια, η διασφάλιση πληροφοριών και η κυβερνοανθεκτικότητα απαντούν σε διαφορετικά ερωτήματα: την αποτροπή της παραβίασης, την εμπιστοσύνη στην πληροφορία και τη συνέχιση της λειτουργίας παρά την αστοχία.
  \item Η τριάδα CIA παραμένει το θεμέλιο, όμως η αυθεντικότητα, η λογοδοσία, η μη αποποίηση, η κατοχή και η χρησιμότητα είναι απαραίτητες για την περιγραφή πολλών πραγματικών καταστάσεων.
  \item Τα μέτρα ασφάλειας επιλέγονται πρώτα με βάση τις αρχές —ελάχιστο προνόμιο, άμυνα σε βάθος, ασφαλείς προεπιλογές, ασφάλεια εκ σχεδιασμού— και μόνο έπειτα με βάση τα προϊόντα.
  \item Ο κίνδυνος προκύπτει όταν μια απειλή συναντά μια ευπάθεια σε ένα πολύτιμο περιουσιακό στοιχείο· ο υπολειπόμενος κίνδυνος πρέπει πάντοτε να μειώνεται, να γίνεται αποδεκτός, να μεταφέρεται ή να αποφεύγεται ρητά.
  \item Τα τυπικά μοντέλα και τα σχήματα αξιολόγησης καθιστούν τους ισχυρισμούς ασφάλειας ακριβείς και ανεξάρτητα επαληθεύσιμους.
\end{itemize}

\section*{Ερωτήσεις ανασκόπησης}
\begin{enumerate}
  \item Εξηγήστε, με δικό σας παράδειγμα, γιατί η κυβερνοανθεκτικότητα παραμένει αναγκαία ακόμη και σε έναν οργανισμό με εξαιρετικά προληπτικά μέτρα.
  \item Περιγράψτε ένα σενάριο στο οποίο η βελτίωση της διαθεσιμότητας αποδυναμώνει την εμπιστευτικότητα. Πώς θα συμφιλιώνατε τους δύο στόχους;
  \item Ποια αρχή των Saltzer και Schroeder παραβιάζεται όταν όλοι οι διαχειριστές μοιράζονται έναν κοινό κωδικό διαχειριστή τομέα; Τεκμηριώστε την απάντησή σας.
  \item Γιατί η διαβάθμιση EAL των Common Criteria δεν αρκεί από μόνη της για να κρίνετε αν ένα προϊόν είναι κατάλληλο για το δικό σας περιβάλλον;
\end{enumerate}
